\chapter{مبرهنة الأعداد الأولية}

\section{نطاق الفصل وحالته}

\noindent\textbf{حالة الفصل:} \textenglish{VERIFIED}. يثبت هذا الفصل الصيغة
المركزية
\[
\psi(x)\sim x,
\]
ثم يستنتج
\[
\vartheta(x)\sim x
\qquad\text{و}\qquad
\pi(x)\sim\frac{x}{\log x}.
\]

بدأت كتابة الفصل بعد إغلاق بوابة
\textenglish{EVIDENCE-FIRST / PRE-AUTHORING}، ثم اجتاز فحوص الجودة وبناء
PDF، والتدقيق المنطقي الداخلي، والتحقق المرجعي للمسار المعتمد. تسجل هذه
الجولات في:
\begin{itemize}
  \item \texttt{docs/CHAPTER\_09\_PRE\_AUTHORING\_AUDIT\_2026-07-20.md}.
  \item \texttt{docs/CHAPTER\_09\_LOGIC\_AUDIT\_2026-07-20.md}.
  \item \texttt{docs/CHAPTER\_09\_BIBLIOGRAPHIC\_VERIFICATION\_2026-07-20.md}.
\end{itemize}
ولا تعني حالة \textenglish{VERIFIED} أن الفصل \textenglish{REVIEWED} أو
\textenglish{RELEASE-READY}؛ فالمراجعة الثانية المستقلة لم تجر بعد.

المسار الأساسي تاوبيري عبر مبرهنة Wiener--Ikehara. أما مسار
Hadamard--de la Vallée Poussin الفعال فيعرض منفصلًا، لأن صيغة بيرون وتحويل
المسار ما زالا دينين علميين معلنين. ويبقى برهان Selberg--Erdős الأولي
مسارًا لاحقًا مستقلًا لا يخلط بالمسار التحليلي.

\subsection{المتطلبات السابقة}

يفترض الفصل:
\begin{itemize}
  \item لغة التقارب والجمع الجزئي من الفصل الثاني.
  \item الاستمرار الميرومورفي ومبدأ الهوية من الفصل الثالث.
  \item دالة فون مانغولت وهويتها الالتفافية من الفصل الرابع.
  \item المشتقة اللوغاريتمية لزيتا في \(\Re(s)>1\) من الفصل الخامس.
  \item قطب زيتا البسيط عند \(s=1\) واستمرارها الميرومورفي من الفصل السادس.
  \item متراجحة ديريشليه المثلثية من الفصل السابع.
\end{itemize}

\subsection{نواتج التعلم}

بعد إتمام الفصل ينبغي أن يستطيع القارئ:
\begin{enumerate}
  \item التمييز بين \(\pi(x)\)، و\(\vartheta(x)\)، و\(\psi(x)\).
  \item إثبات حد تشيبيشيف الخطي باستعمال المعامل الثنائي.
  \item ضبط مساهمة القوى الأولية العليا من دون استعمال مبرهنة الأعداد الأولية.
  \item اشتقاق المتراجحة الموزونة للمشتقة اللوغاريتمية لزيتا.
  \item إثبات عدم وجود أصفار لزيتا على الخط \(\Re(s)=1\) من دون دور منطقي.
  \item مطابقة فروض صيغة Wiener--Ikehara مع معاملات فون مانغولت.
  \item استنتاج \(\psi(x)\sim x\)، ثم \(\vartheta(x)\sim x\)، ثم
  \(\pi(x)\sim x/\log x\).
  \item شرح الفرق بين النتيجة النوعية وحد الخطأ الفعال.
\end{enumerate}

\section{دوال عد الأوليات}

لكل \(x\ge1\) نعرف
\[
\pi(x)=\sum_{p\le x}1,
\qquad
\vartheta(x)=\sum_{p\le x}\log p,
\qquad
\psi(x)=\sum_{n\le x}\Lambda(n).
\]

بما أن \(\Lambda(n)=\log p\) إذا كان \(n=p^m\) قوة أولية، وتساوي صفرًا
في غير ذلك، فإن
\[
\psi(x)
=
\sum_{p^m\le x}\log p
=
\sum_{m\ge1}\vartheta\!\left(x^{1/m}\right).
\]
المجموع الأخير منته، إذ لا يظهر حد بعد \(m>\log_2x\).

الهدف العميق هو إثبات \(\psi(x)\sim x\). أما الانتقال من
\(\vartheta(x)\sim x\) إلى \(\pi(x)\sim x/\log x\) فقد ثبت في
\textenglish{\texttt{ANT-THM-02-04}} بواسطة الجمع الجزئي.

\section{حد تشيبيشيف وضبط القوى الأولية العليا}

\begin{lemma}[حد تشيبيشيف ولمّة القوى الأولية العليا]
\label{lem:chebyshev-and-higher-prime-powers}
\resultid{ANT-LEM-09-02}
\provedhere
عندما \(x\ge2\):
\[
\vartheta(x)\ll x,
\qquad
\psi(x)\ll x,
\]
وكذلك
\[
0\le\psi(x)-\vartheta(x)
\ll\sqrt{x}\log x=o(x).
\]
\end{lemma}

\begin{proof}
ليكن \(n\ge1\) عددًا صحيحًا. كل أولي \(p\) يحقق
\(n<p\le2n\) يقسم المعامل الثنائي \(\binom{2n}{n}\). لذلك
\[
\vartheta(2n)-\vartheta(n)
\le
\log\binom{2n}{n}.
\]
ومن
\[
\binom{2n}{n}
\le
\sum_{j=0}^{2n}\binom{2n}{j}
=2^{2n}
\]
نحصل على
\[
\vartheta(2n)-\vartheta(n)
\le2n\log2.
\]

اختر \(k\) بحيث \(2^k\le x<2^{k+1}\). بالتقسيم إلى فترات ثنائية:
\begin{align*}
\vartheta(x)
&\le\vartheta(2^{k+1})\\
&=
\sum_{j=0}^{k}
\bigl(\vartheta(2^{j+1})-\vartheta(2^j)\bigr)\\
&\le
\sum_{j=0}^{k}2^{j+1}\log2
<4x\log2.
\end{align*}
إذن \(\vartheta(x)\ll x\).

ومن الهوية السابقة ينتج
\begin{align*}
0\le\psi(x)-\vartheta(x)
&=
\sum_{2\le m\le\log_2x}
\vartheta\!\left(x^{1/m}\right)\\
&\ll
\sum_{2\le m\le\log_2x}x^{1/m}\\
&\ll
\sqrt{x}\log x.
\end{align*}
وهذا \(o(x)\). وبضم الحد إلى \(\vartheta(x)\ll x\) نحصل على
\(\psi(x)\ll x\).
\end{proof}

هذه اللمّة مستقلة عن مبرهنة الأعداد الأولية. فلا يجوز استبدالها بتقدير
\(\pi(y)\sim y/\log y\)، لأن ذلك يجعل الاستنتاج دائريًا.

\section{المشتقة اللوغاريتمية وتمثيل ستيلتج}

أثبت الفصل الخامس أنه عندما \(\Re(s)>1\):
\[
-\frac{\zeta'}{\zeta}(s)
=
\sum_{n=1}^{\infty}\frac{\Lambda(n)}{n^s}.
\]

\begin{proposition}[تمثيل المشتقة اللوغاريتمية بواسطة \(\psi\)]
\label{prop:psi-stieltjes-representation}
\resultid{ANT-PROP-09-01}
\provedhere
إذا كان \(\Re(s)>1\)، فإن
\[
-\frac{\zeta'}{\zeta}(s)
=
\int_{1^-}^{\infty}x^{-s}\,d\psi(x)
=
s\int_1^{\infty}\psi(x)x^{-s-1}\,dx.
\]
\end{proposition}

\begin{proof}
الدالة \(\psi\) درجية متصلة من اليمين، وتقفز عند العدد الصحيح \(n\)
بمقدار \(\Lambda(n)\). لذلك
\[
\int_{1^-}^{\infty}x^{-s}\,d\psi(x)
=
\sum_{n=1}^{\infty}\frac{\Lambda(n)}{n^s}
=
-\frac{\zeta'}{\zeta}(s).
\]
التقارب المطلق صحيح في \(\Re(s)>1\). وبالتكامل بالتجزئة على \([1,X]\):
\[
\int_{1^-}^{X}x^{-s}\,d\psi(x)
=
\psi(X)X^{-s}
+s\int_1^X\psi(x)x^{-s-1}\,dx.
\]
من اللمّة السابقة \(\psi(X)\ll X\)، فيؤول الحد الطرفي إلى صفر عندما
\(X\to\infty\)، ويتقارب التكامل لأن \(\Re(s)>1\).
\end{proof}

\section{المتراجحة الموزونة لزيتا}

أثبت الفصل السابع
\[
3+4\cos u+\cos2u=2(1+\cos u)^2\ge0.
\]
ولا نسجل هذه الهوية مرة ثانية؛ بل نسجل تطبيقها التحليلي التالي.

\begin{lemma}[المتراجحة الموزونة للمشتقة اللوغاريتمية]
\label{lem:zeta-weighted-log-derivative-inequality}
\resultid{ANT-LEM-09-01}
\provedhere
لكل \(\sigma>1\) ولكل \(t\in\mathbb R\):
\[
-3\frac{\zeta'}{\zeta}(\sigma)
-4\Re\frac{\zeta'}{\zeta}(\sigma+it)
-\Re\frac{\zeta'}{\zeta}(\sigma+2it)
\ge0.
\]
\end{lemma}

\begin{proof}
من التقارب المطلق لمتسلسلة المشتقة اللوغاريتمية:
\begin{align*}
&-3\frac{\zeta'}{\zeta}(\sigma)
-4\Re\frac{\zeta'}{\zeta}(\sigma+it)
-\Re\frac{\zeta'}{\zeta}(\sigma+2it)\\
&\quad=
\sum_{n=1}^{\infty}
\frac{\Lambda(n)}{n^\sigma}
\left(
3+4\cos(t\log n)+\cos(2t\log n)
\right).
\end{align*}
كل معامل \(\Lambda(n)n^{-\sigma}\) غير سالب، والقوس غير سالب من
\textenglish{\texttt{ANT-LEM-07-01}}، فتنتج الدعوى.
\end{proof}

\section{عدم وجود أصفار على الخط \(\Re(s)=1\)}

\begin{theorem}[مبرهنة Hadamard--de la Vallée Poussin النوعية]
\label{thm:zeta-nonvanishing-on-one-line}
\resultid{ANT-THM-09-01}
\provedhere
لكل \(t\in\mathbb R\setminus\{0\}\):
\[
\zeta(1+it)\ne0.
\]
أما النقطة \(s=1\) فهي قطب بسيط لزيتا، وليست صفرًا.
\end{theorem}

\begin{proof}
افترض عكس المطلوب، ولتكن رتبة الصفر عند \(1+it\) مساوية \(m\ge1\).
ولتكن \(m_2\ge0\) رتبة الصفر عند \(1+2it\)، مع \(m_2=0\) إذا لم تكن
تلك النقطة صفرًا.

من قطب زيتا البسيط عند \(1\):
\[
\frac{\zeta'}{\zeta}(\sigma)
=-\frac1{\sigma-1}+O(1)
\qquad(\sigma\downarrow1).
\]
ومن التحليل المحلي عند الصفر ذي الرتبة \(m\):
\[
\Re\frac{\zeta'}{\zeta}(\sigma+it)
=
\frac{m}{\sigma-1}+O(1),
\]
وكذلك
\[
\Re\frac{\zeta'}{\zeta}(\sigma+2it)
=
\frac{m_2}{\sigma-1}+O(1).
\]
إذن الطرف الأيسر في المتراجحة السابقة يساوي
\[
\frac{3-4m-m_2}{\sigma-1}+O(1).
\]
لكن \(3-4m-m_2\le-1\)، فيصبح الطرف سالبًا عندما تكون
\(\sigma-1\) موجبة وصغيرة، وهذا يناقض المتراجحة.
\end{proof}

\begin{remark}[عدم الدور]
لم يستعمل البرهان مبرهنة الأعداد الأولية، ولا الصيغة الصريحة، ولا منطقة
خالية كمية. اعتمد فقط على المشتقة اللوغاريتمية في \(\Re(s)>1\)، وقطب زيتا
البسيط، والمتراجحة المثلثية غير السالبة.
\end{remark}

\section{مبرهنة Wiener--Ikehara}

نستعمل صيغة خاصة مناسبة للمتسلسلات الديريشلية. هذه النتيجة ليست مثبتة داخل
الموسوعة في هذا الإصدار، ولذلك توسم صراحة بأنها مستشهد بها.

\begin{theorem}[Wiener--Ikehara، صيغة المتسلسلات الديريشلية]
\label{thm:wiener-ikehara-dirichlet-series}
\resultid{ANT-THM-09-02}
\citedresult{\cite[Theorem 1.1، الصفحات 1107--1108]{Korevaar2006}}
لتكن
\[
f(s)=\sum_{n=1}^{\infty}\frac{a_n}{n^s},
\qquad a_n\ge0,
\]
متقاربة في \(\Re(s)>1\)، وضع \(S(x)=\sum_{n\le x}a_n\). افترض أن
\(S(n)=O(n)\)، وأن
\[
f(s)-\frac{A}{s-1}
\]
يمتد تحليليًا أو باستمرار إلى نصف المستوى المغلق \(\Re(s)\ge1\). عندئذ
\[
S(n)\sim An.
\]
\end{theorem}

يرجع الأصل التاريخي إلى Wiener وIkehara
\cite{Wiener1930,Ikehara1931}. واعتمدنا صيغة Korevaar لأنها تصرح بفروض
المتسلسلة الديريشلية وحد المجاميع الجزئية بالصورة التي نحتاجها، وتعرض
تطبيقها على مبرهنة الأعداد الأولية في الصفحتين 1108--1109.

\section{إزالة القطب عند الواحد}

من \textenglish{\texttt{ANT-THM-06-01}} توجد دالة \(h\) هولومورفية قرب
\(1\)، تحقق
\[
\zeta(s)=\frac{h(s)}{s-1},
\qquad h(1)=1.
\]
وبتصغير الجوار تكون \(h\) غير منعدمة. لذلك
\[
-\frac{\zeta'}{\zeta}(s)
=
\frac1{s-1}-\frac{h'}{h}(s),
\]
ومن ثم
\[
-\frac{\zeta'}{\zeta}(s)-\frac1{s-1}
=-\frac{h'}{h}(s)
\]
هولومورفية قرب \(1\).

وعند كل نقطة \(1+it\) مع \(t\ne0\)، تكون زيتا هولومورفية وغير منعدمة من
المبرهنة السابقة، فتكون مشتقتها اللوغاريتمية هولومورفية في جوار النقطة.
إذن الباقي بعد طرح \(1/(s-1)\) يمتد تحليليًا عبر كل نقطة من الخط
\(\Re(s)=1\).

\section{الصيغة المركزية لمبرهنة الأعداد الأولية}

\begin{theorem}[مبرهنة الأعداد الأولية بصيغة فون مانغولت]
\label{thm:prime-number-theorem-psi}
\resultid{ANT-THM-09-03}
\provedhere
عندما \(x\to\infty\):
\[
\psi(x)\sim x.
\]
\end{theorem}

\begin{proof}
طبق مبرهنة Wiener--Ikehara على
\[
a_n=\Lambda(n),
\qquad
f(s)=-\frac{\zeta'}{\zeta}(s).
\]
المعاملات غير سالبة، والسلسلة تتقارب في \(\Re(s)>1\) من الفصل الخامس.
كما أن
\[
S(n)=\sum_{k\le n}\Lambda(k)=\psi(n)\ll n
\]
من لمّة تشيبيشيف. وأثبتنا أن
\[
f(s)-\frac1{s-1}
\]
يمتد تحليليًا عبر الخط \(\Re(s)=1\). لذلك تعطي المبرهنة مع \(A=1\):
\[
\psi(n)\sim n.
\]

ولعدد حقيقي \(x\ge1\):
\[
\psi(x)=\psi(\lfloor x\rfloor).
\]
وبما أن \(\lfloor x\rfloor\sim x\)، نحصل على
\[
\psi(x)
=
\psi(\lfloor x\rfloor)
\sim
\lfloor x\rfloor
\sim x.
\]
\end{proof}

\section{الانتقال إلى \(\vartheta\) و\(\pi\)}

\begin{corollary}[الصيغة المكافئة بدالة تشيبيشيف الأولى]
\label{cor:prime-number-theorem-theta}
\resultid{ANT-COR-09-01}
\provedhere
عندما \(x\to\infty\):
\[
\vartheta(x)\sim x.
\]
\end{corollary}

\begin{proof}
من لمّة القوى الأولية العليا:
\[
0\le\psi(x)-\vartheta(x)=o(x).
\]
ومع \(\psi(x)=x+o(x)\) ينتج
\[
\vartheta(x)
=
\psi(x)-\bigl(\psi(x)-\vartheta(x)\bigr)
=x+o(x).
\]
\end{proof}

\begin{corollary}[مبرهنة الأعداد الأولية]
\label{cor:prime-number-theorem-pi}
\resultid{ANT-COR-09-02}
\provedhere
عندما \(x\to\infty\):
\[
\pi(x)\sim\frac{x}{\log x}.
\]
\end{corollary}

\begin{proof}
هذا تطبيق مباشر لـ\textenglish{\texttt{ANT-THM-02-04}} على النتيجة
\(\vartheta(x)\sim x\).
\end{proof}

\section{ما الذي لم يثبته المسار النوعي؟}

النتيجة السابقة لا تعطي وحدها حد خطأ صريحًا في
\[
\psi(x)-x
\qquad\text{أو}\qquad
\pi(x)-\operatorname{Li}(x).
\]
المسار الكلاسيكي الفعال يحتاج:
\begin{enumerate}
  \item منطقة كمية خالية من الأصفار قرب \(\Re(s)=1\).
  \item تقديرات مناسبة لـ\(\zeta'/\zeta\).
  \item صيغة بيرون مقطوعة أو مملسة.
  \item تحويل المسار وضبط التكاملات على الأضلاع.
  \item استخراج الصيغة الصريحة أو تقدير تكاملي مباشر.
\end{enumerate}

في حالة المستودع الحالية، المنطقة الكلاسيكية الخالية مسجلة
\textenglish{\texttt{CITED}} في \textenglish{\texttt{ANT-THM-06-09}}،
والصيغة الصريحة مسجلة \textenglish{\texttt{CITED}} في
\textenglish{\texttt{ANT-THM-06-08}}، أما بيرون وتحويل المسار فبراهينهما
التفصيلية مؤجلة. لذلك لا ننسب إلى الفصل حد خطأ فعالًا.

\section{المسار الأولي المنفصل}

أثبت Selberg وErdős مبرهنة الأعداد الأولية بطرائق أولية بمعنى أنها لا
تستعمل نظرية الدوال العقدية لزيتا. هذا الوصف لا يعني أن البرهان قصير أو
بسيط؛ فهو يبني هويات وتقديرات مستقلة. سيعالج ذلك المسار لاحقًا بوصفه وحدة
منفصلة، ولا تستعمل نتائجه داخل البرهان التاوبيري الحالي.

\section{خريطة الاعتماد المنطقي}

\begin{enumerate}
  \item حد المعامل الثنائي يعطي \(\vartheta(x)\ll x\).
  \item جمع القوى الأولية يعطي \(\psi(x)\ll x\) و
  \(\psi(x)-\vartheta(x)=o(x)\).
  \item المنتج الأويلري يعطي متسلسلة \(-\zeta'/\zeta\).
  \item المتراجحة المثلثية تعطي المتراجحة الموزونة.
  \item المتراجحة والقطب البسيط يمنعان أصفار الخط \(\Re(s)=1\).
  \item عدم الانعدام يزيل جميع الشذوذ الحدية عدا القطب المعلوم عند \(1\).
  \item Wiener--Ikehara تحول المعلومة التحليلية إلى \(\psi(x)\sim x\).
  \item حد القوى العليا يعطي \(\vartheta(x)\sim x\).
  \item الجمع الجزئي السابق يعطي \(\pi(x)\sim x/\log x\).
\end{enumerate}

لا توجد حافة في هذه الخريطة تعود من مبرهنة الأعداد الأولية إلى فرض سابق.

\section{أخطاء شائعة}

\begin{enumerate}
  \item استعمال عدم انعدام زيتا في \(\Re(s)>1\) كما لو كان يشمل الخط
  \(\Re(s)=1\).
  \item القول إن عدم الانعدام على الخط يعطي وحده مبرهنة الأعداد الأولية من
  دون أداة تاوبيرية أو بديل تحليلي.
  \item تطبيق Wiener--Ikehara من دون التحقق من عدم سلبية المعاملات وحد
  المجاميع الجزئية.
  \item تقدير القوى الأولية العليا باستعمال PNT نفسها.
  \item اعتبار الصيغة الصريحة مثبتة داخل الموسوعة مع أن حالتها
  \textenglish{\texttt{CITED}}.
  \item الخلط بين \(\psi(x)\sim x\) وحد خطأ فعال.
  \item وصف برهان Selberg--Erdős بأنه النسخة نفسها بعد حذف زيتا.
\end{enumerate}

\section{أسئلة تفسير البرهان}

\begin{enumerate}
  \item لماذا تدخل \(\psi\) قبل \(\pi\) في المسار التحليلي؟
  \item أين استعملت عدم سلبية \(\Lambda(n)\)؟
  \item لماذا لا يكفي قطب زيتا عند \(1\) من دون معرفة بقية الخط؟
  \item لماذا يجعل الصفر المفترض عند \(1+it\) المتراجحة الموزونة سالبة؟
  \item ما الدور المختلف لكل من حد تشيبيشيف ومبرهنة Wiener--Ikehara؟
  \item لماذا لا يحمل الانتقال من \(\psi\) إلى \(\vartheta\) دورًا؟
  \item أي أداة إضافية نحتاجها للحصول على حد خطأ فعال؟
\end{enumerate}

\section{تمارين}

\subsection{المستوى الأول}
\begin{enumerate}
  \item أثبت أن كل أولي \(p\) يحقق \(n<p\le2n\) يقسم
  \(\binom{2n}{n}\).
  \item أثبت \(\binom{2n}{n}\le4^n\).
  \item تحقق من الهوية
  \[
  \psi(x)=\sum_{m\ge1}\vartheta(x^{1/m}).
  \]
  \item أثبت أن \(\sqrt{x}\log x=o(x)\).
\end{enumerate}

\subsection{المستوى الثاني}
\begin{enumerate}
  \item اشتق المتراجحة الموزونة حدًا حدًا من متسلسلة
  \(-\zeta'/\zeta\).
  \item اكتب التوسع المحلي لـ\(\zeta'/\zeta\) عند صفر من الرتبة \(m\).
  \item تحقق من إزالة القطب عند \(1\) بكتابة
  \(\zeta(s)=h(s)/(s-1)\).
  \item أثبت صيغة التكامل في
  \textenglish{\texttt{ANT-PROP-09-01}} بالتكامل بالتجزئة لستيلتج.
\end{enumerate}

\subsection{المستوى الثالث}
\begin{enumerate}
  \item بين بدقة أن فروض صيغة Wiener--Ikehara المعتمدة تتحقق عند
  \(a_n=\Lambda(n)\).
  \item قارن بين المسار التاوبيري والمسار القائم على بيرون: ما المعلومة
  الإضافية التي يمكن أن يعطيها الثاني؟
  \item فسر لماذا لا يمكن استنتاج حد خطأ من صيغة تقاربية مجردة.
  \item صمم خريطة اعتماد مستقلة للمسار الأولي من دون استعمال نتائج هذا
  الفصل داخل مقدماته.
\end{enumerate}

\section{المراجع ومسار التوسع}

للبرهانين التاريخيين انظر
\cite{Hadamard1896,DeLaValleePoussin1896}. وللأصل التاوبيري انظر
\cite{Wiener1930,Ikehara1931}. والصيغة الدقيقة المعتمدة وتطبيقها على PNT
موجودان في
\cite[Theorem 1.1 والقسم 2، الصفحات 1107--1109]{Korevaar2006}.

للمعالجات القياسية راجع
\cite{Apostol1976,Davenport2000,MontgomeryVaughan2007,
IwaniecKowalski2004,Tenenbaum2015,Titchmarsh1986}.

المرحلة التالية هي المراجعة الثانية المستقلة، لا الدمج. ولا يرفع الفصل إلى
\textenglish{REVIEWED} قبل استلام حكم مستقل ومعالجة ملاحظاته.

\section{خلاصة الفصل}

جاء الحد الخطي الأولي من المعامل الثنائي، ومنع عدم الانعدام أصفار الحافة
اليمنى لزيتا، ثم حولت Wiener--Ikehara هذه المعلومة التحليلية إلى
\(\psi(x)\sim x\). بعد ذلك أزالت لمّة القوى الأولية العليا الفرق بين
\(\psi\) و\(\vartheta\)، وأكمل الجمع الجزئي الانتقال إلى \(\pi\).

هذا يثبت مبرهنة الأعداد الأولية نوعيًا. أما التقديرات الفعالة فتحتاج أدوات
إضافية لم ندع اكتمالها في هذه النسخة.
